% GENERATED FILE. Edit canonical lessons, then run npm run generate:books.
% canonical-source-hash: fnv1a64:de5b18b30c8da8fe
% canonical-lessons: LA-C375-contingo, LA-C375-exsisto, LA-C375-fortuitus, LA-C375-inopinatus, LA-C375-motus

\chapter{Things That Happen}
\label{ch:la-things-that-happen}
\hlchaptermodality{375}

\begin{chapteropening}
\textbf{By the end of this chapter:} \emph{I can say that something befalls someone or arises, call an event chance or unexpected, and speak of a commotion.}

Complete the last lesson of chapter 375: I can say that something befalls someone or arises, call an event chance or unexpected, and speak of a commotion.
\end{chapteropening}

\section[contingō, contingere]{contingō, contingere — to touch; to happen (to), to fall to one's lot}

\label{lesson:LA-C375-contingo}



\begin{quote}
Before the new one: say the Latin for to hinder, then the Latin for to throw into confusion.
\end{quote}

\subsection*{You'll want to know: contingō, contingere}
\textbf{contingō, contingere} — \textquotedblleft{}to touch; to happen (to), to fall to one's lot\textquotedblright{}.

\begin{quote}
Mihi contigit ut eum vidērem.
\end{quote}

\begin{quote}
\emph{I had the good luck to see him.}
\end{quote}

\emph{contingō, contingere, contigī, contactum}. Built on \emph{tangō}, to touch. \emph{Contingit} + dative, often of good luck. English \emph{contact} and \emph{contingent}.

\begin{cousinweb}
\emph{contingō} is built from \emph{con-} and \emph{tangō} \textquotedblleft{}to touch\textquotedblright{} (EtymDB; Lewis \& Short), the vowel weakened inside the compound: what touches you is what happens to you. English \emph{contact} and \emph{contingent} come from it.
\end{cousinweb}

\begin{grammarlens}[title={What you've built}]
One more word for telling, explaining and describing at B1.
\end{grammarlens}

\subsection*{Your turn}
\begin{itemize}
\raggedright
  \item \emph{Say it:} \emph{contingō, contingere}
  \item \emph{Say it:} \emph{contingō, contingere}, once more
  \item \emph{Say it:} a sentence of your own with \emph{contingō, contingere}
  \item \emph{Recall:} say the Latin for to hinder, then the Latin for to throw into confusion, then say \emph{contingō, contingere} again
\end{itemize}

\begin{tcolorbox}[breakable,colback=teal!4,colframe=teal!35!black,title={Before you move on}]
What does contingō, contingere mean? (\textquotedblleft{}To touch; to happen (to), to fall to one's lot\textquotedblright{}.) Say it once more.
\end{tcolorbox}
\section[exsistō, exsistere]{exsistō, exsistere — to arise, to come forth; to exist}

\label{lesson:LA-C375-exsisto}



\begin{quote}
Before the new one: say the Latin for to throw into confusion, then the Latin for to touch.
\end{quote}

\subsection*{You'll want to know: exsistō, exsistere}
\textbf{exsistō, exsistere} — \textquotedblleft{}to arise, to come forth; to exist\textquotedblright{}.

\begin{quote}
Ex parvīs initiīs magnum bellum exstitit.
\end{quote}

\begin{quote}
\emph{From small beginnings a great war arose.}
\end{quote}

\emph{exsistō, exsistere, exstitī}. \emph{Ex} + \emph{sistō}, to set, to stand: to stand out, to appear. English \emph{exist}.

\begin{cousinweb}
\emph{exsistō} is built from \emph{ex} \textquotedblleft{}out\textquotedblright{} and \emph{sistō} \textquotedblleft{}to stand, to set\textquotedblright{} (Lewis \& Short): to step forth. English \emph{exist} comes from it.
\end{cousinweb}

\begin{grammarlens}[title={What you've built}]
One more word for telling, explaining and describing at B1.
\end{grammarlens}

\subsection*{Your turn}
\begin{itemize}
\raggedright
  \item \emph{Say it:} \emph{exsistō, exsistere}
  \item \emph{Say it:} \emph{exsistō, exsistere}, once more
  \item \emph{Say it:} a sentence of your own with \emph{exsistō, exsistere}
  \item \emph{Recall:} say the Latin for to throw into confusion, then the Latin for to touch, then say \emph{exsistō, exsistere} again
\end{itemize}

\begin{tcolorbox}[breakable,colback=teal!4,colframe=teal!35!black,title={Before you move on}]
What does exsistō, exsistere mean? (\textquotedblleft{}To arise, to come forth; to exist\textquotedblright{}.) Say it once more.
\end{tcolorbox}
\section[fortuītus, fortuīta, fortuītum]{fortuītus, fortuīta, fortuītum — by chance, accidental}

\label{lesson:LA-C375-fortuitus}



\begin{quote}
Before the new one: say the Latin for to touch, then the Latin for to arise.
\end{quote}

\subsection*{You'll want to know: fortuītus, fortuīta, fortuītum}
\textbf{fortuītus, fortuīta, fortuītum} — \textquotedblleft{}by chance, accidental\textquotedblright{}.

\begin{quote}
Nōn fortuītus erat ille cāsus.
\end{quote}

\begin{quote}
\emph{That event was no accident.}
\end{quote}

\emph{fortuītus, fortuīta, fortuītum}, built on \emph{fors}, chance; \emph{fortūna} is its cousin. English \emph{fortuitous}.

\begin{cousinweb}
\emph{fortuītus} is built on \emph{fors} \textquotedblleft{}chance\textquotedblright{} (EtymDB; Lewis \& Short), the same word that gives \emph{fortūna}. English \emph{fortuitous} comes from it.
\end{cousinweb}

\begin{grammarlens}[title={What you've built}]
One more word for telling, explaining and describing at B1.
\end{grammarlens}

\subsection*{Your turn}
\begin{itemize}
\raggedright
  \item \emph{Say it:} \emph{fortuītus, fortuīta, fortuītum}
  \item \emph{Say it:} \emph{fortuītus, fortuīta, fortuītum}, once more
  \item \emph{Say it:} a sentence of your own with \emph{fortuītus, fortuīta, fortuītum}
  \item \emph{Recall:} say the Latin for to touch, then the Latin for to arise, then say \emph{fortuītus, fortuīta, fortuītum} again
\end{itemize}

\begin{tcolorbox}[breakable,colback=teal!4,colframe=teal!35!black,title={Before you move on}]
What does fortuītus, fortuīta, fortuītum mean? (\textquotedblleft{}By chance, accidental\textquotedblright{}.) Say it once more.
\end{tcolorbox}
\section[inopīnātus, inopīnāta, inopīnātum]{inopīnātus, inopīnāta, inopīnātum — unexpected}

\label{lesson:LA-C375-inopinatus}



\begin{quote}
Before the new one: say the Latin for to arise, then the Latin for by chance.
\end{quote}

\subsection*{You'll want to know: inopīnātus, inopīnāta, inopīnātum}
\textbf{inopīnātus, inopīnāta, inopīnātum} — \textquotedblleft{}unexpected\textquotedblright{}.

\begin{quote}
Inopīnātum perīculum nōs terruit.
\end{quote}

\begin{quote}
\emph{An unexpected danger frightened us.}
\end{quote}

\emph{inopīnātus, inopīnāta, inopīnātum}: \emph{in-} not + the participle of \emph{opīnor}: not supposed. \emph{Ex inopīnātō}, unexpectedly.

\begin{cousinweb}
\emph{inopīnātus} is built from \emph{in-} \textquotedblleft{}not\textquotedblright{} and \emph{opīnātus}, the participle of \emph{opīnor} \textquotedblleft{}to suppose\textquotedblright{} (Lewis \& Short): not supposed beforehand.
\end{cousinweb}

\begin{grammarlens}[title={What you've built}]
One more word for telling, explaining and describing at B1.
\end{grammarlens}

\subsection*{Your turn}
\begin{itemize}
\raggedright
  \item \emph{Say it:} \emph{inopīnātus, inopīnāta, inopīnātum}
  \item \emph{Say it:} \emph{inopīnātus, inopīnāta, inopīnātum}, once more
  \item \emph{Say it:} a sentence of your own with \emph{inopīnātus, inopīnāta, inopīnātum}
  \item \emph{Recall:} say the Latin for to arise, then the Latin for by chance, then say \emph{inopīnātus, inopīnāta, inopīnātum} again
\end{itemize}

\begin{tcolorbox}[breakable,colback=teal!4,colframe=teal!35!black,title={Before you move on}]
What does inopīnātus, inopīnāta, inopīnātum mean? (\textquotedblleft{}Unexpected\textquotedblright{}.) Say it once more.
\end{tcolorbox}
\section[mōtus, mōtūs]{mōtus, mōtūs — a movement; a stirring (of the mind); a disturbance}

\label{lesson:LA-C375-motus}



\begin{quote}
Before the new one: say the Latin for by chance, then the Latin for unexpected.
\end{quote}

\subsection*{You'll want to know: mōtus, mōtūs}
\textbf{mōtus, mōtūs} — \textquotedblleft{}a movement; a stirring (of the mind); a disturbance\textquotedblright{}.

\begin{quote}
Magnus mōtus in urbe fuit.
\end{quote}

\begin{quote}
\emph{There was a great commotion in the city.}
\end{quote}

\emph{mōtus, mōtūs} m., from \emph{moveō}. \emph{Mōtus animī}, a stirring of the mind; \emph{mōtus terrae}, an earthquake. English \emph{motion} and \emph{motor} belong to the family.

\begin{cousinweb}
\emph{mōtus} is the noun built on \emph{moveō} \textquotedblleft{}to move\textquotedblright{} (Lewis \& Short), from its participle \emph{mōtus}. English \emph{motion} and \emph{motor} belong to the same family.
\end{cousinweb}

\begin{grammarlens}[title={What you've built}]
One more word for telling, explaining and describing at B1.
\end{grammarlens}

\subsection*{Your turn}
\begin{itemize}
\raggedright
  \item \emph{Say it:} \emph{mōtus, mōtūs}
  \item \emph{Say it:} \emph{mōtus, mōtūs}, once more
  \item \emph{Say it:} a sentence of your own with \emph{mōtus, mōtūs}
  \item \emph{Recall:} say the Latin for by chance, then the Latin for unexpected, then say \emph{mōtus, mōtūs} again
\end{itemize}

\begin{tcolorbox}[breakable,colback=teal!4,colframe=teal!35!black,title={Before you move on}]
What does mōtus, mōtūs mean? (\textquotedblleft{}A movement; a stirring (of the mind); a disturbance\textquotedblright{}.) Say it once more.
\end{tcolorbox}
